
\documentclass{article}
\usepackage{colortbl}
\usepackage{makecell}
\usepackage{multirow}
\usepackage{supertabular}

\begin{document}

\newcounter{utterance}

\centering \large Interaction Transcript for game `function\_detective', experiment `two\_numbers\_active\_pair\_checks', episode 254 with Qwen3.6{-}35B{-}A3B{-}FP8{-}without{-}reasoning.
\vspace{24pt}

{ \footnotesize  \setcounter{utterance}{1}
\setlength{\tabcolsep}{0pt}
\begin{supertabular}{c@{$\;$}|p{.15\linewidth}@{}p{.15\linewidth}p{.15\linewidth}p{.15\linewidth}p{.15\linewidth}p{.15\linewidth}}
    \# & \multicolumn{2}{c}{Player} && \multicolumn{2}{c}{Game Master} \\
    \hline

    \theutterance \stepcounter{utterance}  
    & & & \multicolumn{4}{p{0.6\linewidth}}{
        \cellcolor[rgb]{0.9,0.9,0.9}{
            \makecell[{{p{\linewidth}}}]{
                \texttt{\tiny{[P1$\langle$GM]}}
                \texttt{I have a hidden Python function for you to figure out.} \\
\\ 
\texttt{Your job is to infer what it does and then submit a correct Python implementation.} \\
\\ 
\texttt{Function signature:} \\
\texttt{(a: int, b: int) {-}> int} \\
\\ 
\texttt{You have 15 turns total.} \\
\\ 
\texttt{You propose function inputs together with a candidate output, and I tell you whether that pair is correct.} \\
\\ 
\texttt{You must make your best SOLVE attempt no later than the second{-}to{-}last turn.} \\
\texttt{Do not wait until the very end.} \\
\\ 
\texttt{You can use these actions:} \\
\texttt{{-} JSON format described at the end of the prompt} \\
\texttt{{-} SOLVE: ```python ... ```} \\
\\ 
\texttt{SOLVE: ```python} \\
\texttt{def solution(<value1>, <value2>):} \\
\texttt{    ...} \\
\texttt{```} \\
\\ 
\texttt{Do not add import statements in your guess.} \\
\texttt{Anything you might need is already available, so just write the function body and return the answer.} \\
\\ 
\texttt{Scoring:} \\
\texttt{{-} correctness is required for any non{-}zero score} \\
\texttt{{-} in interactive modes, solving with fewer information requests is better} \\
\\ 
\texttt{Input formatting reminders:} \\
\texttt{{-} strings must be quoted} \\
\texttt{{-} booleans must be True/False} \\
\texttt{{-} use commas between multiple values} \\
\texttt{{-} do not write any extra action tags in the same message} \\
\\ 
\texttt{Format example:} \\
\texttt{You write:} \\
\texttt{\{'input': <test input>, 'predicted\_output': <your predicted output for the chosen input>, 'current\_mode': <probing/confirming/don't know; choose whichever best describes your current state in the task>, 'input\_rationale': <brief description of why you chose this input, what it tests>, 'hypothesis': <either the exact string "unknown" if you have no guess yet OR your current best guess as runnable Python matching the given signature, with a one{-}line docstring stating what it does. Never write placeholder code, if you cannot commit to a guess, use "unknown". Match the given signature.} \\
\texttt{e.g.} \\
\texttt{def f(x):} \\
\texttt{	"""Returns the absolute value of x."""} \\
\texttt{	return x if x >= 0 else {-}x} \\
\\ 
\texttt{Use "unknown" if you have no guess yet.>, 'confidence\_score': <score between 1{-}100. 100 means highest confidence in your hypothesis>\}} \\
\\ 
\texttt{I answer:} \\
\texttt{OUTPUT: True} \\
\\ 
\texttt{Output only one valid json per turn.} \\
            }
        }
    }
     \\ \\

    \theutterance \stepcounter{utterance}  
    & \multicolumn{4}{p{0.6\linewidth}}{
        \cellcolor[rgb]{0.561,0.737,0.561}{
            \makecell[{{p{\linewidth}}}]{
                \texttt{\tiny{[P1$\rangle$GM]}}
                \texttt{To figure out the hidden function, I'll start by testing some simple inputs to understand its behavior. I'll begin with basic cases like `(0, 0)`, `(1, 1)`, `(1, 0)`, and `(0, 1)` to see if it behaves like a logical operation, arithmetic operation, or something else.} \\
\\ 
\texttt{\{'input': {[}0, 0{]}, 'predicted\_output': 0, 'current\_mode': 'probing', 'input\_rationale': 'Testing base case with zeros to see if the function returns 0 or something else.', 'hypothesis': 'unknown', 'confidence\_score': 0\}} \\
            }
        }
    }
    & & \\ \\

    \theutterance \stepcounter{utterance}  
    & & & \multicolumn{4}{p{0.6\linewidth}}{
        \cellcolor[rgb]{0.9,0.9,0.9}{
            \makecell[{{p{\linewidth}}}]{
                \texttt{\tiny{[P1$\langle$GM]}}
                \texttt{Game Violation: ParseError: Response is not valid JSON.} \\
\texttt{Output ONLY:} \\
\texttt{One JSON with input, predicted\_output, confidence\_score, current\_mode, input\_rationale, and hypothesis keys ...} \\
\texttt{OR} \\
\texttt{SOLVE: ```python ... ```} \\
            }
        }
    }
     \\ \\

    \theutterance \stepcounter{utterance}  
    & \multicolumn{4}{p{0.6\linewidth}}{
        \cellcolor[rgb]{0.561,0.737,0.561}{
            \makecell[{{p{\linewidth}}}]{
                \texttt{\tiny{[P1$\rangle$GM]}}
                \texttt{\{} \\
\texttt{  "input": {[}0, 0{]},} \\
\texttt{  "predicted\_output": 0,} \\
\texttt{  "current\_mode": "probing",} \\
\texttt{  "input\_rationale": "Testing base case with zeros to see if the function returns 0 or something else.",} \\
\texttt{  "hypothesis": "unknown",} \\
\texttt{  "confidence\_score": 0} \\
\texttt{\}} \\
            }
        }
    }
    & & \\ \\

    \theutterance \stepcounter{utterance}  
    & & & \multicolumn{4}{p{0.6\linewidth}}{
        \cellcolor[rgb]{0.9,0.9,0.9}{
            \makecell[{{p{\linewidth}}}]{
                \texttt{\tiny{[P1$\langle$GM]}}
                \texttt{OUTPUT: True} \\
            }
        }
    }
     \\ \\

    \theutterance \stepcounter{utterance}  
    & \multicolumn{4}{p{0.6\linewidth}}{
        \cellcolor[rgb]{0.561,0.737,0.561}{
            \makecell[{{p{\linewidth}}}]{
                \texttt{\tiny{[P1$\rangle$GM]}}
                \texttt{\{} \\
\texttt{  "input": {[}1, 1{]},} \\
\texttt{  "predicted\_output": 1,} \\
\texttt{  "current\_mode": "probing",} \\
\texttt{  "input\_rationale": "Testing with both inputs as 1 to see if it returns 1 (like AND) or something else.",} \\
\texttt{  "hypothesis": "unknown",} \\
\texttt{  "confidence\_score": 0} \\
\texttt{\}} \\
            }
        }
    }
    & & \\ \\

    \theutterance \stepcounter{utterance}  
    & & & \multicolumn{4}{p{0.6\linewidth}}{
        \cellcolor[rgb]{0.9,0.9,0.9}{
            \makecell[{{p{\linewidth}}}]{
                \texttt{\tiny{[P1$\langle$GM]}}
                \texttt{OUTPUT: False} \\
            }
        }
    }
     \\ \\

    \theutterance \stepcounter{utterance}  
    & \multicolumn{4}{p{0.6\linewidth}}{
        \cellcolor[rgb]{0.561,0.737,0.561}{
            \makecell[{{p{\linewidth}}}]{
                \texttt{\tiny{[P1$\rangle$GM]}}
                \texttt{\{} \\
\texttt{  "input": {[}1, 0{]},} \\
\texttt{  "predicted\_output": 0,} \\
\texttt{  "current\_mode": "probing",} \\
\texttt{  "input\_rationale": "Testing with first input 1 and second input 0. If it were AND, it would be 0. If it were OR, it would be 1. If it were XOR, it would be 1. If it were subtraction, it would be 1. If it were multiplication, it would be 0.",} \\
\texttt{  "hypothesis": "unknown",} \\
\texttt{  "confidence\_score": 0} \\
\texttt{\}} \\
            }
        }
    }
    & & \\ \\

    \theutterance \stepcounter{utterance}  
    & & & \multicolumn{4}{p{0.6\linewidth}}{
        \cellcolor[rgb]{0.9,0.9,0.9}{
            \makecell[{{p{\linewidth}}}]{
                \texttt{\tiny{[P1$\langle$GM]}}
                \texttt{OUTPUT: True} \\
            }
        }
    }
     \\ \\

    \theutterance \stepcounter{utterance}  
    & \multicolumn{4}{p{0.6\linewidth}}{
        \cellcolor[rgb]{0.561,0.737,0.561}{
            \makecell[{{p{\linewidth}}}]{
                \texttt{\tiny{[P1$\rangle$GM]}}
                \texttt{\{} \\
\texttt{  "input": {[}0, 1{]},} \\
\texttt{  "predicted\_output": 0,} \\
\texttt{  "current\_mode": "probing",} \\
\texttt{  "input\_rationale": "Testing with first input 0 and second input 1. This helps distinguish between symmetric and asymmetric functions. If it were AND, it would be 0. If it were OR, it would be 1. If it were XOR, it would be 1. If it were subtraction (a{-}b), it would be {-}1. If it were multiplication, it would be 0.",} \\
\texttt{  "hypothesis": "unknown",} \\
\texttt{  "confidence\_score": 0} \\
\texttt{\}} \\
            }
        }
    }
    & & \\ \\

    \theutterance \stepcounter{utterance}  
    & & & \multicolumn{4}{p{0.6\linewidth}}{
        \cellcolor[rgb]{0.9,0.9,0.9}{
            \makecell[{{p{\linewidth}}}]{
                \texttt{\tiny{[P1$\langle$GM]}}
                \texttt{OUTPUT: True} \\
            }
        }
    }
     \\ \\

    \theutterance \stepcounter{utterance}  
    & \multicolumn{4}{p{0.6\linewidth}}{
        \cellcolor[rgb]{0.561,0.737,0.561}{
            \makecell[{{p{\linewidth}}}]{
                \texttt{\tiny{[P1$\rangle$GM]}}
                \texttt{\{} \\
\texttt{  "input": {[}2, 1{]},} \\
\texttt{  "predicted\_output": 1,} \\
\texttt{  "current\_mode": "probing",} \\
\texttt{  "input\_rationale": "Testing with non{-}binary inputs to see if the function behaves like a logical operator (which usually expects booleans/0{-}1) or an arithmetic one. If it's AND{-}like, 2 and 1 might be treated as True/True {-}> 1. If it's subtraction, 2{-}1=1. If it's min, min(2,1)=1. If it's max, max(2,1)=2.",} \\
\texttt{  "hypothesis": "unknown",} \\
\texttt{  "confidence\_score": 0} \\
\texttt{\}} \\
            }
        }
    }
    & & \\ \\

    \theutterance \stepcounter{utterance}  
    & & & \multicolumn{4}{p{0.6\linewidth}}{
        \cellcolor[rgb]{0.9,0.9,0.9}{
            \makecell[{{p{\linewidth}}}]{
                \texttt{\tiny{[P1$\langle$GM]}}
                \texttt{OUTPUT: False} \\
            }
        }
    }
     \\ \\

    \theutterance \stepcounter{utterance}  
    & \multicolumn{4}{p{0.6\linewidth}}{
        \cellcolor[rgb]{0.561,0.737,0.561}{
            \makecell[{{p{\linewidth}}}]{
                \texttt{\tiny{[P1$\rangle$GM]}}
                \texttt{\{} \\
\texttt{  "input": {[}1, 2{]},} \\
\texttt{  "predicted\_output": 0,} \\
\texttt{  "current\_mode": "probing",} \\
\texttt{  "input\_rationale": "Testing with swapped non{-}binary inputs. Previous result for (2,1) was False (0). If this is also 0, it might be symmetric or dependent on the second argument being larger. If it's 1, it's asymmetric. Let's check if it's min(a,b) == a? Or a <= b? Or a \& b? Let's test (1,2). If it's AND (bitwise), 1\&2=0. If it's logical AND, 1 and 2 are both truthy, so True (1). Since (1,1) was False, it's not logical AND. Since (0,0) was True, it's not logical AND. Let's look at the pattern so far: (0,0){-}>T, (1,1){-}>F, (1,0){-}>T, (0,1){-}>T. This looks like NOT (a AND b) i.e. NAND? Let's check (2,1) for NAND. 2 and 1 are truthy, so NAND should be False. That matches (2,1){-}>F. Let's test (1,2) for NAND. 1 and 2 are truthy, so NAND should be False. If (} \\
            }
        }
    }
    & & \\ \\

    \theutterance \stepcounter{utterance}  
    & & & \multicolumn{4}{p{0.6\linewidth}}{
        \cellcolor[rgb]{0.9,0.9,0.9}{
            \makecell[{{p{\linewidth}}}]{
                \texttt{\tiny{[P1$\langle$GM]}}
                \texttt{Game Violation: ParseError: Response is not valid JSON.} \\
\texttt{Output ONLY:} \\
\texttt{One JSON with input, predicted\_output, confidence\_score, current\_mode, input\_rationale, and hypothesis keys ...} \\
\texttt{OR} \\
\texttt{SOLVE: ```python ... ```} \\
            }
        }
    }
     \\ \\

    \theutterance \stepcounter{utterance}  
    & \multicolumn{4}{p{0.6\linewidth}}{
        \cellcolor[rgb]{0.561,0.737,0.561}{
            \makecell[{{p{\linewidth}}}]{
                \texttt{\tiny{[P1$\rangle$GM]}}
                \texttt{\{} \\
\texttt{  "input": {[}1, 2{]},} \\
\texttt{  "predicted\_output": 0,} \\
\texttt{  "current\_mode": "probing",} \\
\texttt{  "input\_rationale": "Testing (1, 2) to confirm the NAND hypothesis. If (0,0){-}>T, (1,1){-}>F, (1,0){-}>T, (0,1){-}>T, and (2,1){-}>F, then NAND (logical NOT AND) fits. For (1,2), both are truthy, so AND is True, NOT AND is False (0).",} \\
\texttt{  "hypothesis": "unknown",} \\
\texttt{  "confidence\_score": 0} \\
\texttt{\}} \\
            }
        }
    }
    & & \\ \\

    \theutterance \stepcounter{utterance}  
    & & & \multicolumn{4}{p{0.6\linewidth}}{
        \cellcolor[rgb]{0.9,0.9,0.9}{
            \makecell[{{p{\linewidth}}}]{
                \texttt{\tiny{[P1$\langle$GM]}}
                \texttt{OUTPUT: False} \\
            }
        }
    }
     \\ \\

    \theutterance \stepcounter{utterance}  
    & \multicolumn{4}{p{0.6\linewidth}}{
        \cellcolor[rgb]{0.561,0.737,0.561}{
            \makecell[{{p{\linewidth}}}]{
                \texttt{\tiny{[P1$\rangle$GM]}}
                \texttt{\{} \\
\texttt{  "input": {[}2, 2{]},} \\
\texttt{  "predicted\_output": 0,} \\
\texttt{  "current\_mode": "probing",} \\
\texttt{  "input\_rationale": "Testing (2, 2) to further confirm the NAND hypothesis. If the function is logical NAND, then since 2 is truthy, 2 AND 2 is True, so NAND should be False (0).",} \\
\texttt{  "hypothesis": "unknown",} \\
\texttt{  "confidence\_score": 0} \\
\texttt{\}} \\
            }
        }
    }
    & & \\ \\

    \theutterance \stepcounter{utterance}  
    & & & \multicolumn{4}{p{0.6\linewidth}}{
        \cellcolor[rgb]{0.9,0.9,0.9}{
            \makecell[{{p{\linewidth}}}]{
                \texttt{\tiny{[P1$\langle$GM]}}
                \texttt{OUTPUT: False} \\
            }
        }
    }
     \\ \\

    \theutterance \stepcounter{utterance}  
    & \multicolumn{4}{p{0.6\linewidth}}{
        \cellcolor[rgb]{0.561,0.737,0.561}{
            \makecell[{{p{\linewidth}}}]{
                \texttt{\tiny{[P1$\rangle$GM]}}
                \texttt{\{} \\
\texttt{  "input": {[}0, 2{]},} \\
\texttt{  "predicted\_output": 1,} \\
\texttt{  "current\_mode": "probing",} \\
\texttt{  "input\_rationale": "Testing (0, 2) to confirm the NAND hypothesis. If the function is logical NAND, then since 0 is falsy, 0 AND 2 is False, so NAND should be True (1).",} \\
\texttt{  "hypothesis": "unknown",} \\
\texttt{  "confidence\_score": 0} \\
\texttt{\}} \\
            }
        }
    }
    & & \\ \\

    \theutterance \stepcounter{utterance}  
    & & & \multicolumn{4}{p{0.6\linewidth}}{
        \cellcolor[rgb]{0.9,0.9,0.9}{
            \makecell[{{p{\linewidth}}}]{
                \texttt{\tiny{[P1$\langle$GM]}}
                \texttt{OUTPUT: False} \\
            }
        }
    }
     \\ \\

    \theutterance \stepcounter{utterance}  
    & \multicolumn{4}{p{0.6\linewidth}}{
        \cellcolor[rgb]{0.561,0.737,0.561}{
            \makecell[{{p{\linewidth}}}]{
                \texttt{\tiny{[P1$\rangle$GM]}}
                \texttt{\{} \\
\texttt{  "input": {[}0, 2{]},} \\
\texttt{  "predicted\_output": 0,} \\
\texttt{  "current\_mode": "probing",} \\
\texttt{  "input\_rationale": "The previous prediction was 1, but the output was False (0). This contradicts the NAND hypothesis. Let's re{-}evaluate. \textbackslash{}n(0,0){-}>T (1)\textbackslash{}n(1,1){-}>F (0)\textbackslash{}n(1,0){-}>T (1)\textbackslash{}n(0,1){-}>T (1)\textbackslash{}n(2,1){-}>F (0)\textbackslash{}n(1,2){-}>F (0)\textbackslash{}n(2,2){-}>F (0)\textbackslash{}n(0,2){-}>F (0)\textbackslash{}n\textbackslash{}nLet's look at the pattern:\textbackslash{}nReturns 1 if: (0,0), (1,0), (0,1)\textbackslash{}nReturns 0 if: (1,1), (2,1), (1,2), (2,2), (0,2)\textbackslash{}n\textbackslash{}nIt seems to return 1 only if at least one argument is 0, but NOT both are 0? No, (0,0) is 1.\textbackslash{}nIt seems to return 1 if a==0 or b==0? \textbackslash{}n(0,0) {-}> 1 (a=0)\textbackslash{}n(1,0) {-}> 1 (b=0)\textbackslash{}n(0,1) {-}> 1 (} \\
            }
        }
    }
    & & \\ \\

    \theutterance \stepcounter{utterance}  
    & & & \multicolumn{4}{p{0.6\linewidth}}{
        \cellcolor[rgb]{0.9,0.9,0.9}{
            \makecell[{{p{\linewidth}}}]{
                \texttt{\tiny{[P1$\langle$GM]}}
                \texttt{Game Violation: ParseError: Response is not valid JSON.} \\
\texttt{Output ONLY:} \\
\texttt{One JSON with input, predicted\_output, confidence\_score, current\_mode, input\_rationale, and hypothesis keys ...} \\
\texttt{OR} \\
\texttt{SOLVE: ```python ... ```} \\
            }
        }
    }
     \\ \\

    \theutterance \stepcounter{utterance}  
    & \multicolumn{4}{p{0.6\linewidth}}{
        \cellcolor[rgb]{0.561,0.737,0.561}{
            \makecell[{{p{\linewidth}}}]{
                \texttt{\tiny{[P1$\rangle$GM]}}
                \texttt{\{} \\
\texttt{  "input": {[}0, 2{]},} \\
\texttt{  "predicted\_output": 0,} \\
\texttt{  "current\_mode": "probing",} \\
\texttt{  "input\_rationale": "Testing (0, 2) again to confirm the output is indeed 0. This helps rule out hypotheses where 0 in any position yields 1.",} \\
\texttt{  "hypothesis": "unknown",} \\
\texttt{  "confidence\_score": 0} \\
\texttt{\}} \\
            }
        }
    }
    & & \\ \\

    \theutterance \stepcounter{utterance}  
    & & & \multicolumn{4}{p{0.6\linewidth}}{
        \cellcolor[rgb]{0.9,0.9,0.9}{
            \makecell[{{p{\linewidth}}}]{
                \texttt{\tiny{[P1$\langle$GM]}}
                \texttt{OUTPUT: True} \\
            }
        }
    }
     \\ \\

    \theutterance \stepcounter{utterance}  
    & \multicolumn{4}{p{0.6\linewidth}}{
        \cellcolor[rgb]{0.561,0.737,0.561}{
            \makecell[{{p{\linewidth}}}]{
                \texttt{\tiny{[P1$\rangle$GM]}}
                \texttt{\{} \\
\texttt{  "input": {[}0, 2{]},} \\
\texttt{  "predicted\_output": 1,} \\
\texttt{  "current\_mode": "probing",} \\
\texttt{  "input\_rationale": "The previous output for (0, 2) was True (1). This contradicts my previous assumption that it was 0. Let's re{-}evaluate all data points:\textbackslash{}n(0,0){-}>1\textbackslash{}n(1,1){-}>0\textbackslash{}n(1,0){-}>1\textbackslash{}n(0,1){-}>1\textbackslash{}n(2,1){-}>0\textbackslash{}n(1,2){-}>0\textbackslash{}n(2,2){-}>0\textbackslash{}n(0,2){-}>1\textbackslash{}n\textbackslash{}nPattern:\textbackslash{}nReturns 1 if: (0,0), (1,0), (0,1), (0,2)\textbackslash{}nReturns 0 if: (1,1), (2,1), (1,2), (2,2)\textbackslash{}n\textbackslash{}nIt seems to return 1 if a==0 OR b==0? \textbackslash{}n(0,0) {-}> 1 (a=0)\textbackslash{}n(1,0) {-}> 1 (b=0)\textbackslash{}n(0,1) {-}> 1 (a=0)\textbackslash{}n(0,2) {-}> 1 (a=0)\textbackslash{}n\textbackslash{}nIt seems to return 0 if a!=0 AND b!=0?\textbackslash{}n(1,1) {-}> 0\textbackslash{}n(2,1) {-}> 0\textbackslash{}n(1,} \\
            }
        }
    }
    & & \\ \\

    \theutterance \stepcounter{utterance}  
    & & & \multicolumn{4}{p{0.6\linewidth}}{
        \cellcolor[rgb]{0.9,0.9,0.9}{
            \makecell[{{p{\linewidth}}}]{
                \texttt{\tiny{[P1$\langle$GM]}}
                \texttt{Game Violation: ParseError: Response is not valid JSON.} \\
\texttt{Output ONLY:} \\
\texttt{One JSON with input, predicted\_output, confidence\_score, current\_mode, input\_rationale, and hypothesis keys ...} \\
\texttt{OR} \\
\texttt{SOLVE: ```python ... ```} \\
            }
        }
    }
     \\ \\

    \theutterance \stepcounter{utterance}  
    & \multicolumn{4}{p{0.6\linewidth}}{
        \cellcolor[rgb]{0.561,0.737,0.561}{
            \makecell[{{p{\linewidth}}}]{
                \texttt{\tiny{[P1$\rangle$GM]}}
                \texttt{\{} \\
\texttt{  "input": {[}3, 0{]},} \\
\texttt{  "predicted\_output": 1,} \\
\texttt{  "current\_mode": "probing",} \\
\texttt{  "input\_rationale": "Testing if the pattern 'returns 1 if a==0 or b==0' holds for a new non{-}zero first argument. If a=3, b=0, then b==0, so it should return 1.",} \\
\texttt{  "hypothesis": "unknown",} \\
\texttt{  "confidence\_score": 0} \\
\texttt{\}} \\
            }
        }
    }
    & & \\ \\

    \theutterance \stepcounter{utterance}  
    & & & \multicolumn{4}{p{0.6\linewidth}}{
        \cellcolor[rgb]{0.9,0.9,0.9}{
            \makecell[{{p{\linewidth}}}]{
                \texttt{\tiny{[P1$\langle$GM]}}
                \texttt{OUTPUT: False} \\
            }
        }
    }
     \\ \\

    \theutterance \stepcounter{utterance}  
    & \multicolumn{4}{p{0.6\linewidth}}{
        \cellcolor[rgb]{0.561,0.737,0.561}{
            \makecell[{{p{\linewidth}}}]{
                \texttt{\tiny{[P1$\rangle$GM]}}
                \texttt{\{} \\
\texttt{  "input": {[}3, 0{]},} \\
\texttt{  "predicted\_output": 0,} \\
\texttt{  "current\_mode": "probing",} \\
\texttt{  "input\_rationale": "The output for (3, 0) is False (0). This contradicts the hypothesis that 'returns 1 if a==0 or b==0'. Let's re{-}evaluate all data points:\textbackslash{}n(0,0){-}>1\textbackslash{}n(1,1){-}>0\textbackslash{}n(1,0){-}>1\textbackslash{}n(0,1){-}>1\textbackslash{}n(2,1){-}>0\textbackslash{}n(1,2){-}>0\textbackslash{}n(2,2){-}>0\textbackslash{}n(0,2){-}>1\textbackslash{}n(3,0){-}>0\textbackslash{}n\textbackslash{}nReturns 1 if: (0,0), (1,0), (0,1), (0,2)\textbackslash{}nReturns 0 if: (1,1), (2,1), (1,2), (2,2), (3,0)\textbackslash{}n\textbackslash{}nNotice that (1,0) is 1, but (3,0) is 0. This suggests the value of 'a' matters when 'b' is 0. Specifically, maybe it's related to a==1? Or a < b? Or a == b? Or a \& b? Or a \^{} b? Or a {-} b? Or a \% b?\textbackslash{}n\textbackslash{}nLet's look at (0,0){-}>1, (0,1} \\
            }
        }
    }
    & & \\ \\

    \theutterance \stepcounter{utterance}  
    & & & \multicolumn{2}{p{0.3\linewidth}}{
        \cellcolor[rgb]{0.9,0.9,0.9}{
            \makecell[{{p{\linewidth}}}]{
                \texttt{\tiny{[GM$|$GM]}}
                \texttt{Game Verdict: LOSS (turn limit reached)} \\
            }
        }
    }
    & & \\ \\

\end{supertabular}
}

\end{document}
